\section{Pipeline, Hazards, Forwarding, and Stalls}

\subsection{Why pipelining exists}

A single-cycle CPU waits for every part of an instruction to finish before starting the next instruction. Pipelining improves throughput by overlapping stages from different instructions.

An ideal five-stage pipeline completes one instruction per cycle after filling:

\begin{lstlisting}[language={}]
Cycle 1: I1 IF
Cycle 2: I1 ID   I2 IF
Cycle 3: I1 EX   I2 ID   I3 IF
Cycle 4: I1 MEM  I2 EX   I3 ID   I4 IF
Cycle 5: I1 WB   I2 MEM  I3 EX   I4 ID   I5 IF
\end{lstlisting}

\subsection{Pipeline registers}

Between stages, the CPU stores values in pipeline registers. \texttt{cpu\_pipe.v} has registers with names like:

\begin{lstlisting}[language={}]
IF/ID:  ifid_pc, ifid_instr
ID/EX:  idex_pc, idex_imm, idex_rs1d, idex_rs2d, control bits
EX/MEM: exmem_addr, exmem_store, exmem_result, control bits
MEM/WB: memwb_data, memwb_rd, memwb_reg_write
\end{lstlisting}

These registers preserve each instruction’s state as it moves through the pipeline.

\subsection{Data hazards}

A data hazard occurs when a younger instruction needs a result from an older instruction that has not yet reached write-back.

Example:

\begin{lstlisting}[language={}]
add x5, x1, x2
sub x6, x5, x3
\end{lstlisting}

The \texttt{sub} needs \texttt{x5} before the \texttt{add} writes it back.

\subsection{Forwarding}

Forwarding solves many hazards by taking a value from a later pipeline stage and feeding it directly into EX, without waiting for the register file.

\texttt{cpu\_pipe.v} forwards from:

\begin{itemize}
\item EX/MEM to EX;
\item MEM/WB to EX.
\end{itemize}
The forwarding checks compare destination registers against current source registers and ignore register zero.

\subsection{Load-use stall}

A load result is not available until the memory stage. If the very next instruction needs that loaded value, forwarding cannot provide it soon enough.

Example:

\begin{lstlisting}[language={}]
lw  x5, 0(x1)
add x6, x5, x2
\end{lstlisting}

The pipeline inserts a one-cycle stall. In \texttt{cpu\_pipe.v}, this is detected by checking whether the instruction in EX is a load and its destination matches \texttt{rs1} or \texttt{rs2} of the instruction in ID.

\subsection{Control hazards and flushing}

Branches and jumps decide the next PC. In \texttt{cpu\_pipe.v}, they resolve in EX. If a branch or jump is taken, instructions currently in IF and ID are on the wrong path. The CPU flushes them and redirects the PC to the correct target.

This is called a predict-not-taken pipeline because the fetch unit keeps fetching \texttt{PC+4} until EX proves otherwise.

\bigskip
